\documentclass[11pt]{article}

\usepackage[margin=1in]{geometry}
\usepackage[T1]{fontenc}
\usepackage[utf8]{inputenc}
\usepackage{lmodern}
\usepackage{microtype}
\usepackage{amsmath}
\usepackage{amssymb}
\usepackage{booktabs}
\usepackage{graphicx}
\usepackage{float}
\usepackage{hyperref}
\usepackage{xcolor}
\usepackage{enumitem}

\hypersetup{
    colorlinks=true,
    linkcolor=blue,
    citecolor=blue,
    urlcolor=blue
}

\title{Synthetic Data Generation for Pest Detection}
\author{STA 561 Final Project Team}
\date{\today}

\begin{document}

\maketitle

\begin{abstract}
Collecting labeled pest detection data in real commercial kitchens is difficult because infestations are rare, annotations are expensive, and privacy or sanitation constraints limit data collection. This project studies whether synthetic data can serve as a practical substitute for early-stage detector training. We build an automated pipeline that composes pest assets into kitchen backgrounds, exports bounding-box annotations, and trains a transformer-based object detector using the generated dataset. Our current implementation supports three pest classes: rat, mouse, and cockroach. The pipeline includes kitchen-background collection, synthetic image generation with automatic labels, dataset splitting, DETR training, and evaluation utilities. Although the full video-generation stage and final benchmark results are still in progress, the system already establishes a reproducible framework for analyzing how synthetic scene variation affects pest detection performance. This report describes the motivation, pipeline design, modeling choices, experimental protocol, current status, limitations, and the next steps required to close the gap between image-level synthetic training and the final video-based project objective.
\end{abstract}

\section{Introduction}
Automated pest detection is a useful computer vision problem with direct relevance to food safety, public health, and commercial kitchen operations. Early detection of pests such as rats, mice, and cockroaches can reduce contamination risk and help operators intervene before infestations become severe. In practice, however, training a robust pest detector is constrained by data scarcity. Real images or videos of pests in kitchens are difficult to collect at scale, and even when footage is available, bounding-box annotation is labor-intensive.

Synthetic data offers a practical alternative. If realistic synthetic scenes can be generated with accurate labels, then detector development can proceed without depending entirely on expensive real-world annotation. Synthetic data is especially attractive for rare-event problems, where the positive class is inherently underrepresented. In our setting, synthetic generation also allows direct control over scene variables such as pest class, object scale, background clutter, visibility, and lighting.

The goal of this project is to build an end-to-end pipeline that starts from kitchen imagery, generates labeled synthetic pest scenes, and trains a vision-transformer-based detector to localize and classify pests. Our longer-term assignment target is video-based generation with frame-level labels, but the current repository already supports an image-level synthetic training workflow that serves as the core minimum viable system. The central research question is whether synthetic data can effectively train a pest detector and what aspects of synthetic variation matter most for generalization.

\section{Problem Statement}
We study the following supervised detection problem. Let
\[
x \in \mathbb{R}^{H \times W \times 3}
\]
denote an input kitchen image, or more generally let
\[
x_t \in \mathbb{R}^{H \times W \times 3}, \qquad t=1,\dots,T,
\]
denote the frames of a kitchen video. For each image or frame, the model must identify whether a pest is present, classify its type, and predict a bounding box around it. The three target classes in this project are:

\begin{itemize}[leftmargin=1.5em]
    \item rat,
    \item mouse,
    \item cockroach.
\end{itemize}

Accordingly, the label space is
\[
\mathcal{C}=\{\text{rat},\text{mouse},\text{cockroach}\}.
\]
For a single image $x$, the ground-truth annotation is a finite set
\[
y=\{(c_i,b_i)\}_{i=1}^{N_x},
\]
where $N_x$ is the number of pests visible in the image, $c_i \in \mathcal{C}$ is the class label, and
\[
b_i=(x_{c,i},y_{c,i},w_i,h_i)\in[0,1]^4
\]
is the normalized bounding box for the $i$th pest. Thus, the learning problem is set-valued: the detector must predict an unordered collection of object instances rather than a single class label.

We denote the detector by
\[
f_\theta : \mathbb{R}^{H \times W \times 3} \to \widehat{\mathcal{Y}},
\]
where $\theta$ are model parameters and $\widehat{\mathcal{Y}}$ is the space of predicted labeled boxes. Given a training dataset
\[
\mathcal{D}=\{(x^{(m)},y^{(m)})\}_{m=1}^n,
\]
the standard supervised objective is to learn
\[
\hat{\theta}
=
\arg\min_{\theta}
\frac{1}{n}\sum_{m=1}^{n}
\mathcal{L}\bigl(f_\theta(x^{(m)}),y^{(m)}\bigr),
\]
where $\mathcal{L}$ is a detection loss combining classification and localization error. In our case, this detector is implemented using a vision-transformer-based architecture, specifically DETR.

The practical project requirement is more specific: the final system should take a kitchen image as input, automatically generate synthetic pest videos of length $30$--$60$ seconds with frame-level labels, and use those generated data to train a vision-transformer detector. If a generated video contains $T$ frames, then its annotation may be written as
\[
\mathcal{Y}_{1:T} = (y_1,\dots,y_T),
\]
where each
\[
y_t=\{(c_{t,i},b_{t,i})\}_{i=1}^{N_t}
\]
records the pest instances visible in frame $t$. The ideal final detector should generalize to held-out test videos while maintaining a high true detection rate and a low false positive rate.

From a distributional perspective, let $P_{\mathrm{real}}$ denote the target distribution of real kitchen frames and let $P_{\mathrm{syn}}$ denote the synthetic distribution induced by the generator. The broader research question is whether a detector trained on samples from $P_{\mathrm{syn}}$ can achieve strong performance on test data drawn from $P_{\mathrm{real}}$. In other words, we seek to understand whether synthetic data can serve as an effective surrogate for scarce real labeled pest footage despite the potential domain gap
\[
P_{\mathrm{syn}} \neq P_{\mathrm{real}}.
\]

Our current implementation addresses the core training problem first by building an image-level synthetic data pipeline with automatic bounding-box annotation. Formally, this means we first construct
\[
\mathcal{D}_{\mathrm{syn}}^{\mathrm{img}}=\{(x^{(m)},y^{(m)})\}_{m=1}^n
\]
for static images before extending to video-level data
\[
\mathcal{D}_{\mathrm{syn}}^{\mathrm{vid}}=\{(X^{(m)},\mathcal{Y}^{(m)}_{1:T_m})\}_{m=1}^N.
\]
This is a deliberate systems choice: before scaling to videos, we need a reproducible foundation for asset placement, scene randomization, label correctness, and model training. Establishing the image-level pipeline first therefore provides a controlled setting in which to validate the synthetic generation process, the exported annotations, and the end-to-end detector training procedure.
\section{Related Work}
This project lies at the intersection of three areas: object detection, synthetic data generation, and agricultural or pest-related visual recognition.

First, modern object detection has increasingly shifted toward transformer-based architectures. DETR introduced an end-to-end detection framework that removes many hand-engineered post-processing components and frames detection as direct set prediction \cite{carion2020detr}. This makes it a natural baseline for our work, especially because the course project explicitly requires a vision-transformer-style detector.

Second, synthetic data has become a common strategy when real annotation is scarce or expensive. Synthetic pipelines are widely used in robotics, autonomous driving, and industrial inspection because they provide perfectly aligned labels and controlled variation. The main challenge is domain gap: a detector trained on synthetic imagery may overfit to rendering artifacts or unrealistic scene statistics. Domain randomization is a common mitigation strategy, where nuisance factors such as texture, scale, position, clutter, and lighting are varied aggressively to encourage robustness.

Third, visual pest detection has clear real-world relevance but comparatively limited openly available labeled data, especially for kitchen or sanitation environments. Many applied systems rely on trap images, microscopy, or highly constrained agricultural settings rather than unconstrained indoor scenes. That makes synthetic generation especially valuable in our problem setting.

Our contribution is not a new detection architecture. Instead, we build an integrated, reproducible pipeline for synthetic pest scene generation and transformer-based training, then use it as a testbed for studying whether synthetic data is sufficient for this application.

\section{Method}
\subsection{Pipeline Overview}
The repository implements an automated workflow with the following stages:

\begin{enumerate}[leftmargin=1.5em]
    \item collect kitchen background images,
    \item organize pest assets for each target class,
    \item generate labeled synthetic scenes,
    \item split the resulting detection dataset,
    \item train a DETR-based detector,
    \item evaluate detection performance on held-out data.
\end{enumerate}

The design goal is reproducibility. Each stage is exposed through scripts so that the dataset and model artifacts can be regenerated without manual annotation.

\subsection{Background Collection}
Kitchen backgrounds are collected from Places365. We use the dataset to extract kitchen-related indoor scenes, specifically classes such as \texttt{kitchen} and \texttt{restaurant\_kitchen}. These images act as scene backgrounds onto which pest assets are rendered or composited. Using a large scene repository provides environmental diversity and reduces the chance that the detector memorizes a narrow set of layouts.

\subsection{Pest Assets}
The current pipeline organizes three-dimensional pest assets under class-specific directories:
\begin{itemize}[leftmargin=1.5em]
    \item \texttt{assets/models/rat/},
    \item \texttt{assets/models/mouse/},
    \item \texttt{assets/models/cockroach/}.
\end{itemize}

The use of 3D assets is important because it allows future extension to video generation, viewpoint variation, animation, and more realistic scene interaction. Even in the image-level stage, explicit asset-based generation provides tighter control than generic image editing.

\subsection{Synthetic Scene Generation}
Synthetic training examples are produced with Blender-based automation. We model the generator as a stochastic scene-construction process. For each sample, a kitchen background
\[
B \sim P_B,
\]
a pest asset
\[
A \sim P_A,
\]
and a collection of nuisance parameters
\[
\phi \sim P_{\Phi}
\]
are sampled, where $\phi$ may include object position, scale, in-plane rotation, lighting, occlusion, and other rendering variables. The synthetic image is then generated by a rendering/compositing map
\[
x = G(B,A,\phi),
\]
where $x \in \mathbb{R}^{H \times W \times 3}$ denotes the final RGB image. The corresponding annotation is produced automatically from the scene geometry,
\[
y = T(A,\phi),
\]
where $T$ extracts the pest class and its visible bounding box from the rendered scene.

Equivalently, the synthetic dataset may be viewed as a collection
\[
\mathcal{D}_{\mathrm{syn}}=\{(x_i,y_i)\}_{i=1}^n,
\]
with each pair generated independently from the sampling-and-rendering pipeline above. This formulation makes clear that scene diversity is controlled by the distributions $P_B$, $P_A$, and $P_{\Phi}$: increasing the support of these distributions increases variation in backgrounds, pest appearance, and nuisance factors, which is important for reducing overfitting to narrow synthetic scene statistics.

In practice, for each generated sample the pipeline selects a kitchen background, places a pest instance, applies randomized scene parameters, and exports both the rendered image and the corresponding bounding-box label. Output files include:

\begin{itemize}[leftmargin=1.5em]
    \item a rendered image,
    \item a normalized bounding-box label file,
    \item metadata describing the generated sample.
\end{itemize}

The current generator focuses on image-level examples rather than full video clips. However, the same stochastic formulation extends naturally to temporal data. For a video sequence of length $T$, one may write
\[
x_t = G(B_t,A_t,\phi_t), \qquad t=1,\dots,T,
\]
where the parameters are allowed to evolve over time according to motion paths, camera variation, or time-varying visibility. Thus, the current image-level generator serves as the static special case of a more general framewise synthetic video model.

\subsection{Label Format}
Bounding boxes are exported automatically in normalized coordinates. For a rendered pest instance with pixel-space bounding box
\[
(x_{\min},y_{\min},x_{\max},y_{\max}),
\]
we define the center-width-height representation
\[
x_c=\frac{x_{\min}+x_{\max}}{2}, \qquad
y_c=\frac{y_{\min}+y_{\max}}{2}, \qquad
w=x_{\max}-x_{\min}, \qquad
h=y_{\max}-y_{\min}.
\]
If the rendered image has width $W$ and height $H$, then the normalized box is
\[
\tilde{b}
=
\left(
\frac{x_c}{W},
\frac{y_c}{H},
\frac{w}{W},
\frac{h}{H}
\right)
\in [0,1]^4.
\]
Accordingly, each training annotation can be represented as
\[
y=(c,\tilde{b}),
\]
where $c \in \{\text{rat},\text{mouse},\text{cockroach}\}$ is the class label and $\tilde{b}$ is the normalized bounding box.

Because the scene is generated synthetically, the label is obtained directly from the underlying scene geometry rather than from manual human annotation. In other words, if $\Pi(\cdot)$ denotes projection from scene coordinates to image coordinates, then the exported box is determined by the rendered object support:
\[
\tilde{b} = \operatorname{BBox}\bigl(\Pi(A,\phi)\bigr),
\]
followed by normalization by $(W,H)$. This guarantees consistency between the image and the annotation by construction. One of the main advantages of synthetic data in this context is therefore that annotation correctness can be enforced algorithmically, avoiding the cost and variability of external labeling workflows.
\subsection{Detector}
We use DETR as the baseline detector \cite{carion2020detr}. DETR formulates object detection as a direct set-prediction problem, which removes the need for hand-designed proposal generation or non-maximum suppression and makes it well aligned with the project requirement to use a vision-transformer-style detector for pest localization and classification.

Let an input image be denoted by
\[
x \in \mathbb{R}^{H \times W \times 3}.
\]
After passing $x$ through a convolutional backbone, DETR produces a spatial feature map
\[
F \in \mathbb{R}^{d \times h \times w},
\]
which is flattened into a sequence of visual tokens and processed by a transformer encoder. A transformer decoder then takes a fixed set of learned object queries
\[
\{q_1,\dots,q_Q\}, \qquad q_i \in \mathbb{R}^d,
\]
and outputs $Q$ prediction slots. Each slot predicts a class probability vector and a bounding box:
\[
\hat{y}_i = (\hat{p}_i,\hat{b}_i), \qquad i=1,\dots,Q,
\]
where $\hat{p}_i \in \Delta^{K+1}$ is a probability distribution over $K$ pest classes plus one special ``no-object'' class, and
\[
\hat{b}_i = (\hat{c}_{x,i},\hat{c}_{y,i},\hat{w}_i,\hat{h}_i) \in [0,1]^4
\]
is a normalized bounding box. In our project, the foreground classes are
\[
K=3, \qquad \{\text{rat},\text{mouse},\text{cockroach}\}.
\]

For an image with $N$ ground-truth objects, the target annotation is a set
\[
y = \{(c_j,b_j)\}_{j=1}^N,
\]
where $c_j \in \{1,\dots,K\}$ is the pest class and
\[
b_j = (c_{x,j},c_{y,j},w_j,h_j) \in [0,1]^4
\]
is the normalized ground-truth box. Since DETR always outputs a fixed number $Q$ of predictions with $Q \ge N$, training requires matching the unordered predictions to the unordered ground-truth objects. This is done by bipartite matching via the Hungarian algorithm:
\[
\sigma^\star = \arg\min_{\sigma \in \mathfrak{S}_Q}
\sum_{j=1}^{N}
\mathcal{C}\bigl((c_j,b_j),(\hat{p}_{\sigma(j)},\hat{b}_{\sigma(j)})\bigr),
\]
where $\mathfrak{S}_Q$ denotes the set of permutations of $\{1,\dots,Q\}$ and $\mathcal{C}$ is a matching cost combining classification and localization terms.

After matching, the training loss is applied only to the assigned prediction slots. A standard DETR objective is
\[
\mathcal{L}
=
\sum_{j=1}^{N}
\left[
-\log \hat{p}_{\sigma^\star(j)}(c_j)
+ \lambda_{1}\|b_j-\hat{b}_{\sigma^\star(j)}\|_1
+ \lambda_{\mathrm{giou}}\mathcal{L}_{\mathrm{GIoU}}(b_j,\hat{b}_{\sigma^\star(j)})
\right]
+
\sum_{i \notin \{\sigma^\star(1),\dots,\sigma^\star(N)\}}
-\log \hat{p}_i(\varnothing),
\]
where $\varnothing$ denotes the no-object class, $\|\cdot\|_1$ is the coordinate-wise $\ell_1$ loss, and $\mathcal{L}_{\mathrm{GIoU}}$ is the generalized IoU loss. The first summation encourages correct pest classification and accurate box localization for matched objects, while the second penalizes unmatched queries unless they predict no object.

This formulation is particularly attractive for our setting because each synthetic image contains a small number of pests and comes with automatically generated bounding boxes. Thus, DETR can be trained end-to-end on the synthetic dataset without additional post-processing heuristics. After dataset generation, we split the data into training, validation, and test subsets and optimize the detector parameters by minimizing
\[
\hat{\theta}
=
\arg\min_{\theta}
\frac{1}{n}
\sum_{m=1}^{n}
\mathcal{L}\bigl(f_{\theta}(x^{(m)}), y^{(m)}\bigr),
\]
where $f_{\theta}$ denotes the DETR model and $\{(x^{(m)},y^{(m)})\}_{m=1}^n$ are the synthetic training examples. The repository includes scripts for both training and evaluation, enabling end-to-end experiments once a generated dataset is available.
\section{Experimental Setup}
\subsection{Dataset Construction}
At the time of writing, the repository contains a smoke-test synthetic dataset with 200 labeled images balanced across the three classes. A corresponding split file partitions the data into:

\begin{itemize}[leftmargin=1.5em]
    \item 140 training images,
    \item 30 validation images,
    \item 30 test images.
\end{itemize}

This dataset is not large enough to support strong claims about final model performance, but it is sufficient for pipeline validation, training-script debugging, and early-stage sanity checks.

\subsection{Training Protocol}
The default training configuration uses:
\begin{itemize}[leftmargin=1.5em]
    \item 20 epochs,
    \item batch size 4,
    \item learning rate $1 \times 10^{-4}$.
\end{itemize}

These settings are intended for a baseline run rather than a fully optimized training schedule. As the dataset grows, training hyperparameters and augmentation policies will need to be revisited.

\subsection{Evaluation Metrics}
The final assignment emphasizes true detection rate and false positive rate. To define these quantities precisely, let a predicted bounding box $\hat{b}$ with class label $\hat{c}$ be counted as a true positive if (i) the predicted class matches a ground-truth class $c$, (ii) the prediction confidence exceeds a chosen threshold $\tau_{\mathrm{conf}}$, and (iii) the overlap between prediction and ground truth exceeds a chosen IoU threshold $\tau_{\mathrm{IoU}}$. The intersection-over-union between a predicted box $\hat{b}$ and a ground-truth box $b$ is
\[
\operatorname{IoU}(\hat{b},b)
=
\frac{|\hat{b}\cap b|}{|\hat{b}\cup b|}.
\]
A detection is therefore considered correct when
\[
\hat{c}=c
\qquad\text{and}\qquad
\operatorname{IoU}(\hat{b},b)\ge \tau_{\mathrm{IoU}}.
\]

Given counts of true positives ($TP$), false positives ($FP$), false negatives ($FN$), and true negatives ($TN$), we compute the standard metrics
\[
\operatorname{Precision}=\frac{TP}{TP+FP},
\qquad
\operatorname{Recall}=\frac{TP}{TP+FN}.
\]
In this project, the true detection rate is equivalent to recall:
\[
\operatorname{TDR}=\frac{TP}{TP+FN}.
\]
The false positive rate is
\[
\operatorname{FPR}=\frac{FP}{FP+TN}.
\]
These metrics are useful because they separate two different failure modes: missed pest detections increase $FN$ and reduce TDR, while spurious detections on background clutter increase $FP$ and raise FPR.

For multi-class evaluation, we compute precision and recall separately for each pest class
\[
k \in \{\text{rat},\text{mouse},\text{cockroach}\},
\]
so that
\[
\operatorname{Precision}_k=\frac{TP_k}{TP_k+FP_k},
\qquad
\operatorname{Recall}_k=\frac{TP_k}{TP_k+FN_k}.
\]
Per-class metrics are important because the three pest categories may differ substantially in apparent size, texture, and visibility, which can lead to class-specific detection behavior.

In addition, we report mean average precision (mAP) at a fixed IoU threshold. Let
\[
P_k(r)
\]
denote the precision-recall curve for class $k$ obtained by varying the confidence threshold. The average precision for class $k$ is
\[
\operatorname{AP}_k
=
\int_0^1 P_k(r)\,dr,
\]
or its standard discrete approximation in implementation. The mean average precision is then
\[
\operatorname{mAP}
=
\frac{1}{K}\sum_{k=1}^K \operatorname{AP}_k,
\]
where $K=3$ in our setting. This summary metric is useful because it evaluates ranking quality across confidence thresholds rather than at only one operating point.

The current evaluation script supports thresholded detection analysis with configurable confidence and IoU settings, i.e.,
\[
\tau_{\mathrm{conf}} \in [0,1],
\qquad
\tau_{\mathrm{IoU}} \in [0,1],
\]
so that model behavior can be studied under different operating conditions. This is especially relevant in pest detection, where deployment priorities may favor higher recall or lower false alarm rates depending on the application context.

For the final report, we plan to present:
\begin{itemize}[leftmargin=1.5em]
    \item per-class precision and recall,
    \item overall true detection rate,
    \item false positive rate,
    \item qualitative failure examples.
\end{itemize}

In addition to the aggregate metrics above, qualitative failure examples will help interpret whether errors are caused by small-object scale, occlusion, cluttered backgrounds, or synthetic-to-real mismatch. Together, these quantitative and qualitative measures provide a more complete picture of detector performance than any single scalar metric alone.
\section{Current Results and Project Status}
The engineering pipeline is substantially further along than the final empirical evaluation. At the time of writing, the following components are implemented:

\begin{itemize}[leftmargin=1.5em]
    \item automated Places365 download and kitchen-image filtering,
    \item organization of pest assets by class,
    \item Blender-based synthetic image generation,
    \item automatic bounding-box export,
    \item dataset split generation,
    \item DETR training and evaluation scripts.
\end{itemize}

In addition, a sample synthetic dataset has already been produced locally, and the repository contains the scripts required to regenerate it.

The main missing empirical artifact is a confirmed trained checkpoint and final evaluation report. As a result, we do not yet report final quantitative detector performance. This is an important limitation: the current project status supports a systems and methodology write-up, but not a conclusive answer to the full research question. Rather than fabricating metrics, we leave the main result table in placeholder form below.

\begin{table}[H]
    \centering
    \caption{Planned result table for the final submission.}
    \begin{tabular}{lcccc}
        \toprule
        Training Data & Precision & Recall & TDR & FPR \\
        \midrule
        Synthetic only (smoke dataset) & TBD & TBD & TBD & TBD \\
        Synthetic only (full dataset) & TBD & TBD & TBD & TBD \\
        Synthetic + additional randomization & TBD & TBD & TBD & TBD \\
        Synthetic video frames & TBD & TBD & TBD & TBD \\
        \bottomrule
    \end{tabular}
\end{table}

Even without final metrics, the current pipeline already establishes several useful conclusions. First, the project is technically feasible in a reproducible way. Second, synthetic generation with automatic labels is a practical mechanism for bootstrapping pest detection data. Third, the largest remaining challenges are no longer dataset plumbing or training infrastructure, but realism, scale, and synthetic-to-real generalization.

\section{Discussion}
\subsection{Why Synthetic Data Is Promising Here}
This problem is well suited to synthetic data for several reasons. Pest appearances are relatively structured, the number of target classes is small, and annotation quality is critical. Synthetic generation makes it possible to produce arbitrarily many labeled examples at low marginal cost once the scene pipeline is established. It also enables targeted stress testing: we can deliberately vary clutter, object scale, background type, or visibility to understand failure modes.

\subsection{Limitations of the Current System}
The current system still has several limitations.

First, the synthetic pipeline is image-level rather than fully video-based. This means the repository does not yet satisfy the temporal component of the final assignment. Second, the present dataset scale is modest, so any model trained on it is likely to underfit the intended deployment scenario or overfit to generator artifacts. Third, only a small number of pest assets are currently integrated, which limits appearance diversity. Fourth, the detector has not yet been validated against a realistic held-out video benchmark.

\subsection{Expected Failure Modes}
Several failure modes are likely:
\begin{itemize}[leftmargin=1.5em]
    \item overfitting to synthetic lighting or compositing artifacts,
    \item confusion between small pests and background clutter,
    \item missed detections under heavy occlusion or extreme scale variation,
    \item false positives triggered by kitchen objects with pest-like contours.
\end{itemize}

These failure modes motivate the next phase of work: richer asset variation, improved rendering realism, hard negative examples, and full video generation with realistic motion.

\section{Future Work}
The next development steps are clear.

\begin{enumerate}[leftmargin=1.5em]
    \item Train the baseline DETR model to completion and export a reproducible evaluation report.
    \item Scale the synthetic dataset beyond the current smoke-test size.
    \item Extend the generator from static images to 30--60 second labeled videos.
    \item Add motion policies, occlusion events, and temporal camera variation.
    \item Study synthetic randomization as an ablation variable rather than treating generation as fixed.
    \item Evaluate whether synthetic pretraining plus limited real fine-tuning improves generalization.
\end{enumerate}

From a research perspective, the most interesting question is not simply whether synthetic data works at all, but which forms of synthetic diversity matter most. A carefully designed ablation study could compare background diversity, lighting randomization, object pose, pest size, clutter, and temporal motion to determine which factors most influence detector quality.

\section{Conclusion}
This project develops a reproducible synthetic-data pipeline for pest detection in kitchen environments. The current system collects kitchen backgrounds, integrates pest assets, generates labeled synthetic images, and prepares data for DETR-based training and evaluation. Although the final video-generation stage and quantitative detector results are still pending, the implemented workflow already demonstrates that synthetic data can serve as a practical foundation for this task.

The main contribution of the current project stage is engineering infrastructure: a controllable way to generate labeled pest scenes and train a transformer-based detector without relying on large manually annotated real datasets. The next step is to complete large-scale training and video-based evaluation so that the central research question can be answered with empirical evidence rather than only system design arguments.

\bibliographystyle{plain}
\bibliography{references}

\end{document}
